\chapter{\LaTeX{}使用教程}

\section{字体命令}\label{txt:FreqCmd}
{\kaishu 玲珑骰子安红豆，入骨相思知不知。\hfill ——温庭筠}
	
{\fangsong 愿得一心人，白头不相离。\hfill ——卓文君}
		
{\ifcsname youyuan\endcsname\youyuan\else[无 \cs{youyuan} 字体。]\fi 去年今日此门中，人面桃花相映红。\hfill ——崔护}
			
{\heiti 入我相思门，知我相思苦。\hfill ——李白}
				
{\ifcsname lishu\endcsname\lishu\else[无 \cs{lishu} 字体。]\fi 此情可待成追忆？只是当时已惘然。\hfill ——李商隐}
					
{\songti 雨打梨花深闭门，忘了青春，误了青春。\hfill ——唐寅}

使用\cs{textbf}和\cs{textit}以及\cs{underline}的效果分别如下：

这句话的\textbf{文字}分别\textit{使用}了三种命令来\underline{处理}。

The \textbf{words} in this sentences are \textit{processed} with three different \underline{cmd}.

\section{表格样本}

\subsection{基本表格}
\label{sec:basictable}

模板中关于表格的宏包有三个： \pkg{booktabs}、\pkg{array} 和\pkg{longtabular}。三线表可以用 \pkg{booktabs}提供的 \cs{toprule}、\cs{midrule} 和 \cs{bottomrule}。它们与\pkg{longtable} 能很好的配合使用。
\begin{table}[htb]
	\centering
	\begin{minipage}[t]{0.9\linewidth} % 如果想在表格中使用脚注，minipage是个不错的办法
	\caption[模板文件]{模板文件。如果表格的标题很长，那么在表格索引中就会很不美观，所以要像 chapter 那样在前面用中括号写一个简短的标题。这个标题会出现在索引中。}
	\label{tab:template-files}
	\begin{tabularx}{\linewidth}{lX}
		\toprule
		{\heiti 文件名} & {\heiti 描述} \\
		\midrule
		cquthesis.cls & 模板类文件\footnote{这是一个脚注}\\
		cquthesis.cfg & 模板配置文件\footnote{这是又一个脚注}\\
		cqunumberical.bst & 参考文献 BIB\TeX\ 样式文件。\\
		cquthesis.sty & 常用的包和命令写在这里，减轻主文件的负担。\footnote{同一页上的脚注最多支持到10个}\\
		\bottomrule
		\end{tabularx}
	\end{minipage}
\end{table}

首先来看一个最简单的表格。\ref{tab:template-files} 列举了本模板主要文件及其功能。请大家注意三线表中各条线对应的命令。这个例子还展示了如何在表格中正确使用脚注。由于 \LaTeX{} 本身不支持在表格中使用\cs{footnote}，所以我们不得不将表格放在小页中，而且最好将表格的宽度设置为小页的宽度，这样脚注看起来才更美观。

\begin{table}[htbp]
	\noindent\begin{minipage}{0.5\textwidth}
		\centering
		\caption{第一个并排子表格}
		\label{tab:parallel1}
		\begin{tabular}{p{2cm}p{2cm}}
					\toprule
					No. & Name \\\midrule
					1 & Fox \\
					2 & Panda \\
					3 & Dog \\
					\bottomrule
		\end{tabular}
	\end{minipage}%
	\begin{minipage}{0.5\textwidth}
		\centering
		\caption{第二个并排子表格}{The second subtable in one row}
		\label{tab:parallel2}
		\begin{tabular}{p{2cm}p{2cm}}
			\toprule
			No. & Name \\\midrule
			1 & Charlie \\
			2 & Jack \\
			3 & Tom \\
			\bottomrule
		\end{tabular}
	\end{minipage}
\end{table}

\begin{table}[htbp]
	\centering
	\caption{并排子表格}
	\label{tab:subtable}

	{
		\begin{tabular}{p{2cm}p{2cm}}
			\toprule
			111 & 222 \\\midrule
			222 & 333 \\\bottomrule
		\end{tabular}
	}
	\hskip2cm

	{
		\begin{tabular}{p{2cm}p{2cm}}
			\toprule
			111 & 222 \\\midrule
			222 & 333 \\\bottomrule
		\end{tabular}
	}
\end{table}

不可否认 \LaTeX{} 的表格功能没有想象中的那么强大，不过只要足够认真，足够细致，同样可以排出来非常复杂非常漂亮的表格。

\ref{tab:parallel1}和\ref{tab:parallel2}展示了\cs{xuhao}和\cs{xuhao}\texttt{[1]}的使用，可以达到自动编号的效果。不过要记得在每次使用之前使用\cs{resetxuhao}，或者\cs{xuhao}\texttt{[1]}。使用\cs{setxuhao}\oarg{1-6}可以更改序号的标记方式，如\ref{tab:parallel2}所示。详细用法请参阅用户手册。

\begin{longtable}[c]{c*{6}{r}}
	\caption[实验数据]{实验数据，这个题注是双语的，而且十分的长，注意这在索引中的处理方式}\label{tab:performance}\\
	\toprule
	测试程序 & \multicolumn{1}{c}{正常运行} & \multicolumn{1}{c}{同步} & \multicolumn{1}{c}{检查点} & \multicolumn{1}{c}{卷回恢复}
	& \multicolumn{1}{c}{进程迁移} & \multicolumn{1}{c}{检查点} \\
	& \multicolumn{1}{c}{时间 (s)}& \multicolumn{1}{c}{时间 (s)}&
	\multicolumn{1}{c}{时间 (s)}& \multicolumn{1}{c}{时间 (s)}& \multicolumn{1}{c}{
		时间 (s)}&  文件（KB）\\\midrule
	\endfirsthead
	\multicolumn{7}{c}{续表~\thetable\hskip1em 实验数据}\\
	\toprule
	测试程序 & \multicolumn{1}{c}{正常运行} & \multicolumn{1}{c}{同步} & \multicolumn{1}{c}{检查点} & \multicolumn{1}{c}{卷回恢复}
	& \multicolumn{1}{c}{进程迁移} & \multicolumn{1}{c}{检查点} \\
	& \multicolumn{1}{c}{时间 (s)}& \multicolumn{1}{c}{时间 (s)}&
	\multicolumn{1}{c}{时间 (s)}& \multicolumn{1}{c}{时间 (s)}& \multicolumn{1}{c}{
		时间 (s)}&  文件（KB）\\\midrule
	\endhead
	\hline
	\multicolumn{7}{r}{续下页}
	\endfoot
	\endlastfoot
	CG.A.2 & 23.05 & 0.002 & 0.116 & 0.035 & 0.589 & 32491 \\
	CG.A.4 & 15.06 & 0.003 & 0.067 & 0.021 & 0.351 & 18211 \\
	CG.A.8 & 13.38 & 0.004 & 0.072 & 0.023 & 0.210 & 9890 \\
	CG.B.2 & 867.45 & 0.002 & 0.864 & 0.232 & 3.256 & 228562 \\
	CG.B.4 & 501.61 & 0.003 & 0.438 & 0.136 & 2.075 & 123862 \\
	CG.B.8 & 384.65 & 0.004 & 0.457 & 0.108 & 1.235 & 63777 \\
	MG.A.2 & 112.27 & 0.002 & 0.846 & 0.237 & 3.930 & 236473 \\
	MG.A.4 & 59.84 & 0.003 & 0.442 & 0.128 & 2.070 & 123875 \\
	MG.A.8 & 31.38 & 0.003 & 0.476 & 0.114 & 1.041 & 60627 \\
	MG.B.2 & 526.28 & 0.002 & 0.821 & 0.238 & 4.176 & 236635 \\
	MG.B.4 & 280.11 & 0.003 & 0.432 & 0.130 & 1.706 & 123793 \\
	MG.B.8 & 148.29 & 0.003 & 0.442 & 0.116 & 0.893 & 60600 \\
	LU.A.2 & 2116.54 & 0.002 & 0.110 & 0.030 & 0.532 & 28754 \\
	LU.A.4 & 1102.50 & 0.002 & 0.069 & 0.017 & 0.255 & 14915 \\
	LU.A.8 & 574.47 & 0.003 & 0.067 & 0.016 & 0.192 & 8655 \\
	LU.B.2 & 9712.87 & 0.002 & 0.357 & 0.104 & 1.734 & 101975 \\
	LU.B.4 & 4757.80 & 0.003 & 0.190 & 0.056 & 0.808 & 53522 \\
	LU.B.8 & 2444.05 & 0.004 & 0.222 & 0.057 & 0.548 & 30134 \\
	CG.B.2 & 867.45 & 0.002 & 0.864 & 0.232 & 3.256 & 228562 \\
	CG.B.4 & 501.61 & 0.003 & 0.438 & 0.136 & 2.075 & 123862 \\
	CG.B.8 & 384.65 & 0.004 & 0.457 & 0.108 & 1.235 & 63777 \\
	MG.A.2 & 112.27 & 0.002 & 0.846 & 0.237 & 3.930 & 236473 \\
	MG.A.4 & 59.84 & 0.003 & 0.442 & 0.128 & 2.070 & 123875 \\
	MG.A.8 & 31.38 & 0.003 & 0.476 & 0.114 & 1.041 & 60627 \\
	MG.B.2 & 526.28 & 0.002 & 0.821 & 0.238 & 4.176 & 236635 \\
	MG.B.4 & 280.11 & 0.003 & 0.432 & 0.130 & 1.706 & 123793 \\
	MG.B.8 & 148.29 & 0.003 & 0.442 & 0.116 & 0.893 & 60600 \\
	LU.A.2 & 2116.54 & 0.002 & 0.110 & 0.030 & 0.532 & 28754 \\
	LU.A.4 & 1102.50 & 0.002 & 0.069 & 0.017 & 0.255 & 14915 \\
	LU.A.8 & 574.47 & 0.003 & 0.067 & 0.016 & 0.192 & 8655 \\
	LU.B.2 & 9712.87 & 0.002 & 0.357 & 0.104 & 1.734 & 101975 \\
	LU.B.4 & 4757.80 & 0.003 & 0.190 & 0.056 & 0.808 & 53522 \\
	LU.B.8 & 2444.05 & 0.004 & 0.222 & 0.057 & 0.548 & 30134 \\
	EP.A.2 & 123.81 & 0.002 & 0.010 & 0.003 & 0.074 & 1834 \\
	EP.A.4 & 61.92 & 0.003 & 0.011 & 0.004 & 0.073 & 1743 \\
	EP.A.8 & 31.06 & 0.004 & 0.017 & 0.005 & 0.073 & 1661 \\
	EP.B.2 & 495.49 & 0.001 & 0.009 & 0.003 & 0.196 & 2011 \\
	EP.B.4 & 247.69 & 0.002 & 0.012 & 0.004 & 0.122 & 1663 \\
	EP.B.8 & 126.74 & 0.003 & 0.017 & 0.005 & 0.083 & 1656 \\
	\bottomrule
\end{longtable}

如果你要排版的表格长度超过一页，那么推荐使用 \pkg{longtable} 或者 \pkg{supertabular}宏包，模板对 \pkg{longtable} 进行了相应的设置，所以用起来可能简单一些。表~\ref{tab:performance} 就是 \pkg{longtable} 的简单示例。


\section{数学公式}
\label{sec:equation}
贝叶斯公式如式~(\ref{equ:chap1:bayes})，其中 $p(y|\mathbf{x})$ 为后验；
$p(\mathbf{x})$ 为先验；分母 $p(\mathbf{x})$ 为归一化因子。
\begin{equation}
\label{equ:chap1:bayes}
p(y|\mathbf{x}) = \frac{p(\mathbf{x},y)}{p(\mathbf{x})}=
\frac{p(\mathbf{x}|y)p(y)}{p(\mathbf{x})} 
\end{equation}

论文里面公式越多，\TeX{} 就越 happy。再看一个 \pkg{amsmath} 的例子：
\newcommand{\envert}[1]{\left\lvert#1\right\rvert} 
\begin{equation}\label{detK2}
\det\mathbf{K}(t=1,t_1,\dots,t_n)=\sum_{I\in\mathbf{n}}(-1)^{\envert{I}}
\prod_{i\in I}t_i\prod_{j\in I}(D_j+\lambda_jt_j)\det\mathbf{A}
^{(\lambda)}(\overline{I}|\overline{I})=0.
\end{equation} 

前面定理示例部分列举了很多公式环境，可以说把常见的情况都覆盖了，大家在写公式的时候一定要好好看 \pkg{amsmath} 的文档，并参考模板中的用法：
\begin{multline*}%\tag{[b]} % 这个出现在索引中的
\int_a^b\biggl\{\int_a^b[f(x)^2g(y)^2+f(y)^2g(x)^2]
-2f(x)g(x)f(y)g(y)\,dx\biggr\}\,dy \\
=\int_a^b\biggl\{g(y)^2\int_a^bf^2+f(y)^2
\int_a^b g^2-2f(y)g(y)\int_a^b fg\biggr\}\,dy
\end{multline*}

这里还有一个多级规划公式，这个公式使用\csgo{listeq}{索引名}手动加入了目录后的索引。
\begin{equation}\label{bilevel}
\left\{\begin{array}{l}
\max\limits_{{\mbox{\footnotesize\boldmath $x$}}} F(x,y_1^*,y_2^*,\cdots,y_m^*)\\[0.2cm]
\mbox{subject to:}\\[0.1cm]
\qquad G(x)\le 0\\[0.1cm]
\qquad(y_1^*,y_2^*,\cdots,y_m^*)\mbox{ solves problems }(i=1,2,\cdots,m)\\[0.1cm]
\qquad\left\{\begin{array}{l}
\max\limits_{{\mbox{\footnotesize\boldmath $y_i$}}}f_i(x,y_1,y_2,\cdots,y_m)\\[0.2cm]
\mbox{subject to:}\\[0.1cm]
\qquad g_i(x,y_1,y_2,\cdots,y_m)\le 0.
\end{array}\right.
\end{array}\right.
\end{equation}
这些跟规划相关的公式都来自于清华大学刘宝碇老师《不确定规划》的课件。以上的许多例子由清华大学的薛瑞尼同学编写。

\section{化学方程式}

使用\pkg{mhchem}的\csgo{ce}{化学式或方程式}能够让你很容易地表示出各种化学式和化学方程：

例如：
\begin{center}
	\ce{C6H5-CHO}\\ \ce{A\bond{~--}B\bond{~=}C\bond{-~-}D}\\ \ce{SO4^2- + Ba^2+ -> BaSO4 v}
\end{center}

复杂一点的方程式也不在话下，如\eqref{eq:chem}：
\begin{equation}\label{eq:chem}
	\ce{Zn^2+
		<=>[+ 2OH-][+ 2H+]
		$\underset{\text{amphoteres Hydroxid}}{\ce{Zn(OH)2 v}}$ <=>[+ 2OH-][+ 2H+]
		$\underset{\text{Hydroxozikat}}{\ce{[Zn(OH)4]^2-}}$
	}
\end{equation}

这个方程式嵌套在了\pkg{equation}环境中，可用\cs{eqlist}(\cs{listeq}的别名，作用相同)来编排到索引中。


\section{国际单位制(SI Unit)}

模板采用\pkg{siunitx}作为国际单位制支持宏包，以下是一些使用例子，这个包的文档写得非常不错，请在命令行里输入\texttt{texdoc siunitx}察看。
\begin{center}
	\num{.3e45}\\
	\num{1.654 x 2.34 x 3.430}\\
	\si{\kilogram\metre\per\second}\\    
	\SIlist{0.13;0.67;0.80}{\milli\metre}
\end{center}


\section{绘图}
\label{sec:draw}

本模板不预先装载任何绘图包（如 \pkg{pstricks，pgf} 等），完全由用户来决定。个人觉得 \pkg{pgf} 不错，不依赖于 Postscript。此外还有很多针对 \LaTeX{} 的GUI 作图工具，如 XFig(jFig), WinFig, Tpx, Ipe, Dia, Inkscape, LaTeXPiX,jPicEdt, jaxdraw 等等。

\section{插图}
\label{sec:graphs}

推荐《\LaTeXe\ 插图指南》。关于子图形的使用细节请参看 \pkg{subcaption} 宏包的说明文档。

\subsection{一个图形}
\label{sec:onefig}
一般图形都是处在浮动环境中。之所以称为浮动是指最终排版效果图形的位置不一定与源文
件中的位置对应\footnote{这是\LaTeX 的一个设计特性。}，这也是刚使
用 \LaTeX{} 同学可能遇到的问题。如果要强制固定浮动图形的位置，请使用 \pkg{float} 宏包，
它提供了 \texttt{[H]} 参数，比如图~\ref{fig:xfig1}。
\begin{figure}[htb] % use float package if you want it here
	\centering
	\includegraphics[height=4cm]{bitlogo.pdf}
	\caption{北京理工大学大学校徽}
	\label{fig:xfig1}
\end{figure}

大学之道，在明明德，在亲民，在止于至善。知止而后有定；定而后能静；静而后能安；安
而后能虑；虑而后能得。物有本末，事有终始。知所先后，则近道矣。古之欲明明德于天
下者，先治其国；欲治其国者，先齐其家；欲齐其家者，先修其身；欲修其身者，先正其心；
欲正其心者，先诚其意；欲诚其意者，先致其知；致知在格物。物格而后知至；知至而后
意诚；意诚而后心正；心正而后身 修；身修而后家齐；家齐而后国治；国治而后天下
平。自天子以至于庶人，壹是皆以修身为本。其本乱而未治者 否矣。其所厚者薄，而其所
薄者厚，未之有也！

\hfill —— 《大学》


\subsection{多个图形}
\label{sec:multifig}

如果多个图形相互独立，并不共用一个图形计数器，那么用 \texttt{minipage} 或者\texttt{parbox} 就可以。否则，请参看
\begin{figure}[ht]
	\centering%
	\begin{subfigure}{3cm}
		\includegraphics[height=4cm]{bitlogo.pdf}
		\caption{第一个小图形}
	\end{subfigure}%
	\hspace{4em}%
	\begin{subfigure}{3cm}
		\includegraphics[height=3cm]{bitlogo.pdf}
		\caption{第二个小图形，注意这个图略矮些。subfigure中同一行的子图在顶端对齐。}
	\end{subfigure}
	\caption{包含子图形的大图形(subfigure示例)}
	\label{fig:big1-subfigure}
\end{figure}

如果要把编号的两个图形并排，那么小页就非常有用了。
\begin{figure}
	\begin{minipage}{0.48\textwidth}
		\centering
		\includegraphics[height=5cm]{bitlogo.pdf}
		\caption{并排第一个图}
		\label{fig:parallel1}
	\end{minipage}\hfill
	\begin{minipage}{0.48\textwidth}
		\centering
		\includegraphics[height=5cm]{bitlogo.pdf}
		\caption{并排第二个图}
		\label{fig:parallel2}
	\end{minipage}
\end{figure}

测试用途：theequation值为：\theequation ，thefigure值为：\thefigure ，thetable值为：\thetable
